\section{Representation and physical-measurement follow-ups}
\label{app:physics-followup}
The new experiments extend the historical controls with narrower questions: whether future-state prediction preserves useful expansion increments, whether a scalar motion quantity is observable under the fitted model, and whether a sensor or clock residual predicts a better TTC estimate. Each experiment retains its own population and reference model. An improvement in an auxiliary metric does not substitute for a TTC comparison. \evidence{N04,N05}

\subsection{A closed predictive-learning pilot}
The event student learns against spatial features from a frozen, distilled V-JEPA~2.1 ViT-B teacher~\citep{vjepa21}. For each of seeds 7, 13, and 21, the pilot runs static distillation, one-step prediction, and a cross-sequence shuffled-future control for 1,000 updates with batch size 16. TTC labels do not enter these student losses. The final checkpoint is fixed before evaluation. After training, a ridge probe with regularization parameter one reads temporal differences of the frozen encoder, with the predictor removed. It fits on 512 P samples from twenty-four sequences and evaluates 256 H samples from eight different sequences. Box-based localization remains privileged information. \evidence{N04}

The primary target is annotated log-height change. Its error is dimensionless: it is neither a time-normalized rate nor TTC error in seconds. The criterion requires the one-step arm to beat both static distillation and shuffled futures, with the upper endpoint of each sequence-bootstrap difference interval below zero. The bootstrap resamples eight complete sequences while keeping the three seeds together, so its interval remains conditional on those seeds and on historically exposed H sequences.

\EvidenceJepaTable
The candidate's probe MAE is 0.03564004, compared with 0.03763239 for static distillation and 0.03506189 for shuffled futures. Its candidate-minus-static difference is $-0.00199235$, with interval $[-0.00496210,0.00081898]$; candidate-minus-shuffled is $0.00057816$, with interval $[-0.00217005,0.00292590]$. Neither superiority criterion passes. The random-initialization control gives 0.03719641. This failed continuation criterion is not an equivalence test or evidence that predictive learning can never help.

The auxiliary latent metrics clarify the limitation. Future teacher-state MSE improves relative to both controls, but teacher-increment MSE worsens. The corresponding candidate-minus-static changes are $-0.01719163$ for state and $+0.03848016$ for increment; candidate-minus-shuffled changes are $-0.05317388$ and $+0.00447422$. Predicting a future state can compensate for present-state errors or model variation irrelevant to expansion, but the experiment does not identify either explanation causally. Its results establish that lower state-prediction error was insufficient for the prescribed physical probe. The pilot closed without rollout, scaling, or integration into the released estimator. \evidence{N04}

\subsection{Target-specific observability and model error}
For a linearized observation model $\mathbf y=H\boldsymbol\beta+\boldsymbol\epsilon$ and target $q=\mathbf c^\top\boldsymbol\beta$, noiseless identifiability requires $\mathbf c$ to lie in the row space of $H$. In the affine image-motion parameterization, isotropic expansion uses $\mathbf c=(1/2,0,0,1/2,0,0)^\top$. The target can be estimable without every affine coefficient being identifiable. Conversely, a pseudoinverse always returning a number does not certify that the target is supported. This algebraic condition is classical and is used here as a diagnostic. \evidence{N05}

On the 149-query screen, the diagnostic marks $q$ estimable in 63 queries. Among those estimates, 26 have an incorrect sign, ten lie outside the canonical phase domain, and eight have zero residual degrees of freedom. A small fitted residual in the last group cannot demonstrate reliability. The target itself also need not equal local isotropic image expansion when rotation, parallax, or changing visible surfaces violate the measurement model.

The gain-prediction experiment compares physical descriptors against simpler support descriptors with leave-one-acquisition-out fitting. Physical-feature MSE is 0.919267453, versus 0.492043720 for support features; their difference is 0.427223732 with interval $[0.080550690,1.010417232]$. This is MSE for predicting correction gain, not TTC MAE. The result rejects the proposed descriptor gate. A historical robust readout returns only eight usable estimates, four with the wrong sign; after preserving F0 fallback, its MAE is 0.779543461\,s versus 0.766024827\,s for the matched R0 baseline. The comparison records unavailable measurements and the original missing output; the error of the eight supported measurements is not presented as complete-population performance. \evidence{N05}

\subsection{Sensor ambiguity and the missing reference}
The sensor-identifiability bank constructs 24 pairs of matching event streams, with pixel, polarity, and timestamp agreement to a tolerance of one picosecond. Twelve pairs have different inverse TTC. Under the bank's log-affine texture model, changing motion and inversely scaling the texture gradient can preserve the increments that generate events; the distinct initial references are absent from the event stream. Providing a known reference helps distinguish the restricted conditions. This construction diagnoses an ambiguity in the tested model, not a general impossibility theorem for event-based estimation. \evidence{N05}

The real-data proxy uses 72 road queries from twenty-four acquisitions. Slider is excluded from the photometric fit because its alignment remains provisional. On common support, observed RGB change gives sensor-residual MSE 96.657290 versus 145.648563 for the physical proxy. Its improvement over an activity control is less decisive: the corresponding descriptive difference interval includes zero. Sensor-residual reductions cannot be interpreted as TTC improvements, and the inventory does not establish calibrated contrast thresholds, gain, initial reference, or refractory response.

The direct rate readout tests that implication. The affine photometric model permits translation, deformation, and global illumination variation and returns 72 valid rate estimates. With the original fallback retained, its macro TTC MAE is 9.168523738\,s, compared with 0.766024827\,s for R0. The spatially permuted control is worse still, at 45.629594312\,s. Beating that control is insufficient when the matched estimator deteriorates so strongly. The experiment opens no producer or JEPA training and supplies no calibrated confidence signal for fusion. \evidence{N05}

\subsection{Surfaces and dynamic clocks}
For positive depth $Z$, signed inverse TTC is $q=-\dot Z/Z$. Differentiation gives
\begin{equation}
 \dot q=q^2-\ddot Z/Z.
\end{equation}
If axial relative velocity is constant in this coordinate, its finite-time evolution is
\begin{equation}
 q(t+\delta)=\frac{q(t)}{1-\delta q(t)},\qquad T(t+\delta)=T(t)-\delta.
\end{equation}
The transport cannot cross the contact singularity. Acceleration, camera rotation,
or a change in the visible surface can invalidate the assumptions even when an
object has constant velocity in world coordinates. These relations follow from
kinematics; the experiment tests their utility in the estimator rather than
claiming a new physical law. \evidence{N05}

A three-frame surface route supplies 133 of 149 measurements and falls back on the other sixteen. Its system MAE is 3.103982972\,s versus R0's 0.766024827\,s, with a paired difference interval $[0.934231309,4.198244944]$\,s. The crop audit finds no changing crop within the 146 auditable histories; 53 contain a constant anisotropic resize. This rules out one proposed changing-zoom explanation for that cohort, not errors from occlusion, backgrounds, or parallax. \evidence{N05}

The constant-axial-velocity transport law above is exact only under its model assumptions. Reapplying its composition identity inside an exact implementation adds no new measurement. The follow-up compares fixed persistence, smoothing, transport, and observer rules against their matched event context after canonical PAIR. Context macro MAE is 1.941890262\,s; EWMA gives 1.952822072\,s, the observer 1.981365164\,s, and persistence/transport 3.104231759\,s. These comparisons do not establish a benefit. Their baseline is the event context, which must remain distinct from the multimodal R0 baseline used above. \evidence{N05}

Together, these experiments distinguish implementation correctness, model identifiability, and useful downstream prediction. No rejected gate becomes a successful learned physical model because its equations are interpretable. Further training would require a new, bounded hypothesis supported by these failures and evaluated with appropriate controls; it is not a result already achieved by this revision.
\FloatBarrier
